\documentclass[a4paper,12pt]{article}
\usepackage[T2A]{fontenc}
\usepackage[utf8]{inputenc}
\usepackage[russian]{babel}
\usepackage{amsmath,amssymb,mathtools}
\usepackage{helvet}
\renewcommand{\familydefault}{\sfdefault}
\usepackage{icomma}
\usepackage[left=20mm,right=20mm,top=20mm,bottom=22mm]{geometry}
\usepackage{microtype}
\usepackage{tikz}
\usetikzlibrary{calc,arrows.meta,positioning,shapes.geometric}
\usepackage{tkz-euclide}
\usepackage{pgfplots}
\pgfplotsset{compat=1.18}
\usepackage{tabularx,booktabs}
\usepackage{enumitem}
\usepackage{parskip}
\usepackage{fancyhdr}
\usepackage{setspace}
\usepackage{xcolor}
\definecolor{accent}{HTML}{2A6F73}
\definecolor{darkaccent}{HTML}{17484B}
\definecolor{textgray}{HTML}{333333}
\definecolor{softgray}{HTML}{E8EEEE}
\definecolor{warning}{HTML}{8A4B35}
\color{textgray}
\onehalfspacing
\setlength{\parindent}{0pt}
\setlength{\headheight}{25pt}
\setlist[itemize]{leftmargin=6mm,itemsep=2pt,topsep=3pt}
\setlist[enumerate]{leftmargin=8mm,itemsep=4pt,topsep=4pt}
\pagestyle{fancy}
\fancyhf{}
\fancyhead[L]{\small\color{darkaccent}Чек-лист по математике}
\fancyhead[R]{\small\color{darkaccent}27.09.26}
\fancyfoot[C]{\small\thepage}
\renewcommand{\headrulewidth}{0.4pt}
\newcommand{\keyidea}[1]{\medskip{\color{darkaccent}\bfseries #1}\smallskip}
\newcommand{\important}[1]{{\color{warning}\bfseries #1}}

\begin{document}

\begin{titlepage}
\thispagestyle{empty}
\vspace*{24mm}
\begin{center}
{\Large\color{darkaccent}\bfseries ЧЕК-ЛИСТ}

\vspace{9mm}
{\huge\bfseries Сложные задачи ОГЭ:\\[3mm]
неравенства, графики, движение\\
и подобие треугольников}

\vspace{12mm}
{\large Повторение ключевых методов занятия}

\vspace{22mm}
\begin{tabular}{rl}
\textbf{Ученица:} & София Кальней \\[3mm]
\textbf{Преподаватель:} & Артём Александрович Лёвин \\[3mm]
\textbf{Дата:} & 27.09.26
\end{tabular}

\vfill
{\large\color{darkaccent}
Лёвин Артём Александрович\\
эксклюзивно для Софии Кальней}

\vspace{18mm}
\begin{tikzpicture}[x=1cm,y=1cm]
\draw[accent,line width=1.4pt]
(0,0).. controls (3,0.8) and (6,-0.8) .. (9,0)
.. controls (11.5,0.8) and (14,-0.4) .. (16.4,0.25);
\end{tikzpicture}
\end{center}
\end{titlepage}

\section*{\color{darkaccent}1. Сложные неравенства: сначала ищите структуру}

Перед механическим раскрытием скобок проверьте, можно ли преобразовать выражение короче.

\keyidea{Разность квадратов}

Если
\[
A^2\ge B^2,
\]
удобно перенести всё в одну сторону:
\[
A^2-B^2\ge0,
\qquad
A^2-B^2=(A-B)(A+B).
\]

\subsection*{Пример}
\[
(9x-4)^2\ge(4x-9)^2.
\]
\[
(9x-4)^2-(4x-9)^2\ge0,
\]
\[
(5x+5)(13x-13)\ge0,
\]
\[
65(x+1)(x-1)\ge0.
\]
Ответ:
\[
\boxed{x\in(-\infty;-1]\cup[1;+\infty)}.
\]

\keyidea{Почему разность квадратов особенно полезна при больших степенях}

\[
(2x^2-1)^2\le(x^2-2)^2.
\]
\[
(2x^2-1)^2-(x^2-2)^2\le0,
\]
\[
3(x^2+1)(x^2-1)\le0.
\]
Поскольку
\[
x^2+1>0,
\]
остаётся
\[
x^2-1\le0.
\]
Ответ:
\[
\boxed{x\in[-1;1]}.
\]

\keyidea{Оценка знака множителя}

\[
x^4+x^{100}+47>0
\]
при любом действительном \(x\).

\important{Нельзя делить неравенство на выражение с переменной, пока не установлен его знак и невозможность обращения в нуль.}

\subsection*{Чётные степени}
\[
a^{2n}\ge0,
\qquad
(-a)^{2n}=a^{2n}.
\]
Например,
\[
(3x-1)^4=(1-3x)^4.
\]

\section*{\color{darkaccent}2. Сумма неотрицательных выражений}

Если
\[
A\ge0,\qquad B\ge0,
\]
то
\[
A+B=0
\]
только при
\[
\boxed{A=0\quad\text{и}\quad B=0}.
\]

Это работает для квадратов, модулей, чётных степеней и арифметических корней чётной степени.

\subsection*{Пример}
\[
|x^2-x-6|+|x^2+4x+4|=0.
\]
Необходимо одновременно:
\[
\begin{cases}
x^2-x-6=0,\\
x^2+4x+4=0.
\end{cases}
\]
\[
x^2-x-6=(x-3)(x+2),
\qquad
x^2+4x+4=(x+2)^2.
\]
Общий корень:
\[
\boxed{x=-2}.
\]

\section*{\color{darkaccent}3. Средняя скорость}

\[
\boxed{v_{\text{ср}}=\frac{S_{\text{общ}}}{t_{\text{общ}}}}.
\]

Для равных половин пути:
\[
\boxed{
v_{\text{ср}}=\frac{2v_1v_2}{v_1+v_2}
}.
\]

Для \(66\) км/ч и \(78\) км/ч:
\[
v_{\text{ср}}
=
\frac{2\cdot66\cdot78}{66+78}
=
\boxed{71{,}5\text{ км/ч}}.
\]

\section*{\color{darkaccent}4. Многоэтажные дроби}

\[
\boxed{\frac{\frac ab}{\frac cd}=\frac{ad}{bc}},
\qquad
\boxed{\frac{\frac ab}{c}=\frac{a}{bc}},
\qquad
\boxed{\frac{a}{\frac bc}=\frac{ac}{b}}.
\]

\section*{\color{darkaccent}5. Рациональная функция и выколотая точка}

Рассмотрим
\[
y=\frac{(x-3)(x^2+6x+8)}{x+2}.
\]
ОДЗ:
\[
x\ne-2.
\]
Так как
\[
x^2+6x+8=(x+2)(x+4),
\]
при \(x\ne-2\)
\[
y=(x-3)(x+4)=x^2+x-12.
\]

Выколотая точка:
\[
y(-2)=-10,
\qquad
\boxed{(-2;-10)}.
\]

Вершина:
\[
x_{\text{в}}=-\frac12,
\qquad
y_{\text{в}}=-\frac{49}{4}.
\]

Для горизонтали \(y=m\) ровно одна общая точка получается при
\[
\boxed{
m=-\frac{49}{4}
\quad\text{или}\quad
m=-10
}.
\]

\section*{\color{darkaccent}6. Подобие треугольников}

Пусть
\[
AB=10,\quad BC=14,\quad BF=5,\quad BE=7,\quad AC=15.
\]
Тогда
\[
\frac{BF}{BA}=\frac{5}{10}=\frac12,
\qquad
\frac{BE}{BC}=\frac{7}{14}=\frac12,
\]
а угол при \(B\) общий. Поэтому
\[
\triangle BFE\sim\triangle BAC,
\]
и
\[
EF=\frac12\cdot15=\boxed{7{,}5}.
\]

В трапеции \(ABCD\), где \(AD\parallel BC\), диагонали пересекаются в точке \(O\):
\[
\frac{AO}{OC}=\frac{AD}{BC}.
\]
Если
\[
AD=16,\quad BC=9,\quad AC=15,
\]
то
\[
\frac{x}{15-x}=\frac{16}{9},
\qquad
\boxed{AO=9{,}6}.
\]

\section*{\color{darkaccent}7. Самопроверка}

\begin{itemize}
\item[$\square$] Я замечаю разность квадратов до раскрытия скобок.
\item[$\square$] Я умею оценивать знак множителя.
\item[$\square$] Я не делю неравенство на выражение неизвестного знака.
\item[$\square$] Я понимаю, когда сумма неотрицательных выражений равна нулю.
\item[$\square$] Я использую полный путь и полное время при поиске средней скорости.
\item[$\square$] Я сначала выписываю ОДЗ рациональной функции.
\item[$\square$] После сокращения возвращаю выколотую точку.
\item[$\square$] В геометрии ищу общий или равный угол и пропорции сходственных сторон.
\end{itemize}

\clearpage
\section*{\color{darkaccent}Тренировочный блок --- Путь А}

\begin{enumerate}
\item Решите:
\[
(5x-2)^2\ge(2x-5)^2.
\]

\item Решите:
\[
(3x^2-2)^2\le(x^2-4)^2.
\]

\item Решите:
\[
(x^8+x^4+9)(2x-5)\le0.
\]

\item Решите:
\[
|x^2-5x+6|+|x^2-4x+4|=0.
\]

\item Решите:
\[
(2x-1)^4\ge(2x+1)^4.
\]

\item Первую половину пути автомобиль прошёл со скоростью \(60\) км/ч, вторую --- со скоростью \(90\) км/ч. Найдите среднюю скорость.

\item Вычислите:
\[
\frac{\frac56}{\frac{10}{9}}.
\]

\item Дана функция
\[
y=\frac{(x-1)(x^2+5x+6)}{x+2}.
\]
Найдите все \(m\), при которых прямая \(y=m\) имеет с графиком ровно одну общую точку.

\item В треугольнике \(ABC\)
\[
AB=15,\quad BC=10,\quad AC=20,
\]
\(BF=6\), \(BE=4\). Найдите \(EF\).

\item В трапеции \(ABCD\), \(AD\parallel BC\), диагонали пересекаются в \(O\). Известно:
\[
AD=15,\quad BC=10,\quad AC=25.
\]
Найдите \(AO\).
\end{enumerate}

\clearpage
\section*{\color{darkaccent}Ответы}

\renewcommand{\arraystretch}{1.5}
\begin{center}
\begin{tabularx}{\linewidth}{@{}cX@{}}
\toprule
\textbf{№} & \textbf{Ответ} \\
\midrule
1 & \(x\in(-\infty;-1]\cup[1;+\infty)\) \\
2 & \(x\in[-\sqrt{3/2};\sqrt{3/2}]\) \\
3 & \(x\in(-\infty;5/2]\) \\
4 & \(x=2\) \\
5 & \(x\le0\) \\
6 & \(72\text{ км/ч}\) \\
7 & \(3/4\) \\
8 & \(m=-4,\,-3\) \\
9 & \(EF=8\) \\
10 & \(AO=15\) \\
\bottomrule
\end{tabularx}
\end{center}

\vfill
\begin{center}
{\small\color{darkaccent}
Лёвин Артём Александрович эксклюзивно для Софии Кальней}
\end{center}

\end{document}
